\section{Finding Patterns}
\label{sec:cses-finding-patterns}

\problemstatement{Given a text $t$ and $k$ patterns, determine for each
pattern whether it appears as a substring of $t$. Output ``YES'' or
``NO'' per pattern.}

\begin{patternbox}
\emph{``Do these many patterns occur in one fixed text?''} is the
signature use of a \textbf{suffix automaton (SAM)} built on $t$. The SAM
recognises exactly the substrings of $t$, so each pattern is a matter of
walking its characters through the automaton.
\end{patternbox}

\subsection*{The suffix automaton recognises every substring}

A suffix automaton of $t$ is the smallest automaton whose paths from the
initial state spell precisely the substrings of $t$. Building it takes
$O(|t|)$ (with a constant or $\log\Sigma$ factor for transitions). To test
a pattern $p$, start at the initial state and follow the transition for
each character of $p$; if at any point the needed transition is missing,
$p$ is not a substring, otherwise it is. The total query cost is the sum
of pattern lengths.

\begin{ideabox}[Build Once, Query Many]
The SAM front-loads all the work into a single linear construction over
$t$; afterwards every pattern is answered in time proportional to its own
length. This ``preprocess the text, stream the queries'' shape is why SAM
dominates multi-pattern substring questions.
\end{ideabox}

\subsection*{Algorithm}

\begin{enumerate}
  \item Build the suffix automaton of $t$ by extending it one character at
        a time.
  \item For each pattern, walk from the initial state following its
        characters; report YES if the walk completes, else NO.
\end{enumerate}

\subsection*{Dry run}

\dryrun{$t=$``aybabtu''; patterns ``bab'', ``ba'', ``xyz''.}

\begin{itemize}
  \item ``bab'': walk b$\to$a$\to$b succeeds -- YES.
  \item ``ba'': succeeds -- YES.
  \item ``xyz'': the transition on ``x'' is missing at the start -- NO.
\end{itemize}

\subsection*{C++ implementation}

\begin{lstlisting}
#include <bits/stdc++.h>
using namespace std;

struct SAM {
    struct State { int len, link; map<char,int> next; };
    vector<State> st;
    int last;
    SAM() { st.push_back({0, -1, {}}); last = 0; }
    void extend(char c) {
        int cur = st.size();
        st.push_back({st[last].len + 1, -1, {}});
        int p = last;
        while (p != -1 && !st[p].next.count(c)) { st[p].next[c] = cur; p = st[p].link; }
        if (p == -1) st[cur].link = 0;
        else {
            int q = st[p].next[c];
            if (st[p].len + 1 == st[q].len) st[cur].link = q;
            else {
                int clone = st.size();
                st.push_back(st[q]);
                st[clone].len = st[p].len + 1;
                while (p != -1 && st[p].next[c] == q) { st[p].next[c] = clone; p = st[p].link; }
                st[q].link = st[cur].link = clone;
            }
        }
        last = cur;
    }
};

int main() {
    string t;
    int k;
    cin >> t >> k;
    SAM sam;
    for (char c : t) sam.extend(c);

    while (k--) {
        string p;
        cin >> p;
        int cur = 0;
        bool ok = true;
        for (char c : p) {
            auto it = sam.st[cur].next.find(c);
            if (it == sam.st[cur].next.end()) { ok = false; break; }
            cur = it->second;
        }
        cout << (ok ? "YES" : "NO") << "\n";
    }
    return 0;
}
\end{lstlisting}

\begin{complexitybox}
\textbf{Time Complexity.} $O(|t|\log\Sigma)$ to build, plus
$O(\sum|p_i|\log\Sigma)$ for all queries. \textbf{Space Complexity.}
$O(|t|)$ states.
\end{complexitybox}

\begin{takeawaybox}
A suffix automaton turns ``is this a substring?'' into a walk, so many
patterns against a fixed text are answered by one construction plus cheap
traversals. The same SAM, augmented with counts, answers the next two
problems (Counting Patterns, Pattern Positions) with almost no extra code.
\end{takeawaybox}
